% Generated draft: requires mathematical and source review.
\documentclass[11pt]{article}
\usepackage[margin=25mm]{geometry}
\usepackage{fontspec}
\setmainfont{Latin Modern Roman}
\usepackage{amsmath,amssymb,mathtools,unicode-math}
\setmathfont{Latin Modern Math}
\usepackage{enumitem,needspace,xurl,hyperref}
\setlength{\parindent}{0pt}
\setlength{\parskip}{5pt}
\setlength{\emergencystretch}{3em}
\newcommand{\N}{\mathbb{N}}
\newcommand{\Z}{\mathbb{Z}}
\newcommand{\Q}{\mathbb{Q}}
\newcommand{\R}{\mathbb{R}}
\newcommand{\C}{\mathbb{C}}
\newcommand{\F}{\mathbb{F}}
\newcommand{\A}{\mathbb{A}}
\newcommand{\PP}{\mathbb{P}}
\newcommand{\E}{\mathbb{E}}
\newcommand{\eps}{\varepsilon}
\newcommand{\dd}{\,\mathrm{d}}
\begin{document}
\section*{UnsolvedMath ID: 30003422}
\textbf{Source problem code:} catalogue entry
\subsection{Definitions and mathematical statement}
A differential graded (dg) category is a category whose morphism spaces are cochain complexes over a fixed coefficient base, with composition satisfying the graded Leibniz rule. For a dg category \(\mathcal B\) with a single object \(b\), its endomorphism complex \(\operatorname{End}_{\mathcal B}(b)\) is a dg algebra. Its perfect derived category \(\operatorname{per}(\mathcal B)\) is the smallest full triangulated subcategory of the derived category of right dg \(\mathcal B\)-modules that contains the representable module \(\mathcal B(-,b)\) and is closed under direct summands.

The supplied record uses \(\mathcal B_m\) and \(\mathcal C^{(d)}_{\infty,m}\), calling the latter a higher cluster category, but supplies no definitions of these objects, coefficient base, or admissible ranges for \(d\) and \(m\). Subject to those missing specifications, the requested construction has the following conditional form: for each intended admissible pair \((d,m)\), construct a natural dg category \(\mathcal B^{(d)}_m\) with exactly one object \(b\), together with identifications
\[
\operatorname{End}_{\mathcal B^{(d)}_m}(b)\cong\mathcal B_m,
\qquad
\operatorname{per}\!\left(\mathcal B^{(d)}_m\right)\simeq\mathcal C^{(d)}_{\infty,m}.
\]
The second identification expresses the intended equality of categories as an equivalence compatible with their triangulated structures, if these are the intended structures. The source does not specify the required notion of identification in either expression. In particular, it does not say whether \(\mathcal B_m\) denotes the full dg endomorphism algebra, an underlying graded algebra, or a cohomology algebra. No choice among these interpretations is made here.

A topological interpretation of the construction is additionally desired, rather than stated as a necessary condition. The terms ``natural'' and ``nice,'' and the intended topological interpretation, have no specified mathematical criteria in the record. Consequently, the displayed target is a conditional statement template requiring these source definitions and qualifications before it becomes a fully determined existence problem.
\subsection{Short English statement}
The question asks for a natural one-object dg model of the intended higher cluster category, with prescribed endomorphism algebra, preferably admitting a topological interpretation. The algebra, target category, parameter ranges, coefficient base, and required identifications are unspecified, so the supplied record does not determine a complete mathematical existence statement.
\subsection{Sources}
\begin{itemize}
\item Supplied local UnsolvedMath registry record, permanent ID \(30003422\), title \emph{DG Models for Higher Cluster Categories}, exact source problem code \texttt{catalogue entry}, in the Oberwolfach Reports Open Problems collection. The record gives proposed year \(2017\) and names no proposer.
\item Supplied report citation: \emph{Representation Theory of Quivers and Finite Dimensional Algebras} (2018), Oberwolfach Reports, pp.~591--681, DOI \texttt{10.4171/owr/2017/12}. Supplied source link: \url{https://doi.org/10.4171/owr/2017/12}. The supplied background describes Oberwolfach Reports as documenting mathematical workshops and their open-problem discussions. No precise problem page or individual contribution is supplied.
\item Supplied literature source, ``OWR Representation Theory report'': \url{https://ems.press/content/serial-article-files/46676}.
\item Supplied literature source, ``Higher Nakayama algebras'': \url{https://arxiv.org/abs/1604.03500}.
\item Literature assessment dated \texttt{2026-08-21}: supplied dated curation, not independently verified here. It labels the problem open and states that the report presents the dg-category and topological-realization request as a question, and that higher Nakayama and topological Fukaya developments do not establish the exact requested construction for all parameters. The background describes this as dated literature triage derived from the supplied OWR curation and requests independent verification of time-sensitive status and cited claims.
\item The registry labels the statement \texttt{corrected\_verified} and reports manual semantic, typographic, and LaTeX review confirming self-containment and preservation of the target. These are supplied curation claims; this draft does not independently verify them or the supplied open label.
\end{itemize}
\end{document}
